\documentclass[tikz, border=8pt]{standalone}
\usepackage{tikz}
\usepackage{graphicx}
\usetikzlibrary{calc}

\begin{document}
\begin{tikzpicture}
  %% ── 원본 이미지를 배경으로 삽입 ──────────────────────────────────────
  %% Overleaf에 이미지를 "diagram.png" 이름으로 업로드하세요.
  \node[anchor=south west, inner sep=0] (img) at (0,0)
    {\includegraphics[width=14cm]{보이스 피싱 탐지 모델 구조도}};
  %% 이미지 TikZ 좌표계: 너비 14cm, 높이 4.54cm (비율 2413:783 유지)

  %% ── Local Encoder (파란색) ────────────────────────────────────────────
  %% 커버 영역: Sliding Window + ×T Segments 점선 박스 (RoBERTa + Attention Pooling)
  \draw[blue!65, very thick, rounded corners=10pt,
        fill=blue!10, fill opacity=0.35]
    (1.35, 3.05) rectangle (11.05, 4.50);
  \node[blue!70, font=\bfseries\small, anchor=south west]
    at (1.35, 4.50) {Local Encoder};

  %% ── Global Context Model (초록색) ────────────────────────────────────
  %% 커버 영역: Aux Head + Mamba SSM Block + Last Valid State Gather
  %%           + Normal/Super/Phishing Head + Fusion & Log + Final Output
  \draw[green!50!black, very thick, rounded corners=10pt,
        fill=green!10, fill opacity=0.35]
    (0.08, 0.08) rectangle (13.05, 2.90);
  \node[green!45!black, font=\bfseries\small, anchor=south west]
    at (0.08, 2.90) {Global Context Model};

\end{tikzpicture}
\end{document}
